% !TEX program = xelatex
% “华为杯”第二十三届中国研究生数学建模竞赛论文
% 模板：paper/gmcmthesis.cls（按官方附件3 .doc 逐项实测复刻；规范见 docs/附件2）
% 编译：xelatex main && bibtex main && xelatex main && xelatex main
\documentclass[colorprint]{gmcmthesis}

\baominghao{2026XXXXX}        % 参赛队号（提交前务必核对）
\schoolname{（学校名称）}
\membera{（队员一姓名）}
\memberb{（队员二姓名）}
\memberc{（队员三姓名）}
\title{复杂场景下多模态情感识别的数学建模与算法设计}

\begin{document}
\maketitle                    % 封面(页脚居中"0") + 摘要页页首(页码从 1 起，无页眉)

%% ============ 摘要页 ============
%% 规范要点：题目/摘要/关键词写在摘要页；篇幅一般不超过两页；无需英文摘要；
%% 摘要需含 建模思路、主要方法、模型、结果与结论、创新点、关键词
\begin{abstract}
（摘要正文：针对……问题，本文建立……模型。首先……；其次……；最后……。
主要结果与创新点：……。）

\keywords{多模态情感识别；跨模态时序对齐；缺失鲁棒预测；可解释性}
\end{abstract}

%% ============ 目录（可选） ============
\maketoc

%% ============ 正文（章节对应 paper/outline.md） ============
\section{问题重述}
%（素材：题面 docs/复杂场景下多模态情感识别的数学建模与算法设计.docx）

\section{问题分析与技术路线}
%（素材：P1→P2→P3 闭环，方案 v2 §7 Figure 1；总架构图 figures/ 待出）

\section{模型假设与符号说明}
%（素材：方案 v2 §2/§4：半开区间、20ms 量化、BERT 截断规则等入假设）

\section{问题一：多模态时序对齐模型的建立与求解}
%（素材：docs/问题一建模方案.md v2，含全部实测背书数字）

\section{问题二：面向缺失场景的鲁棒情感预测模型}
%（素材：docs/问题二建模方案.md；Baseline→Enhanced→Innovation 三层、掩码融合、双输出）

\section{问题三：模型可解释性分析}
%（素材：docs/问题三建模方案.md；反事实/贡献度：沙普利、蒸馏、局部缺失）

\section{实验与结果分析}
%（素材：docs/实验假设与验证目标.md；消融矩阵、缺失率×位置×模态设计；表格入 tables/）

\section{模型评价与推广}

%% ============ 参考文献 ============
%% 规范（附件2）：
%%   书籍：[编号] 作者，书名，出版地：出版社，起止页码，出版年。
%%   期刊：[编号] 作者，论文名，杂志名，卷期号：起止页码，出版年。
%%   网络：[编号] 作者，资源标题，网址，访问时间（年月日）。
%% 正文引用处用方括号标示编号，如 [1][3]；引用书籍须指明页码。
\bibliographystyle{gmcm}
\bibliography{reference}

%% ============ 附录 ============
\appendix
\section{环境与核心参数}
%（素材：docs/环境与复现说明.md、environment-manifest.json 字段表）

\section{关键代码}
\begin{Python}{align.py（示例：P1 词区间对齐）}
# TODO: 从 src/ 摘取关键片段
\end{Python}

\end{document}
